\chapter{Experimental Protocol}
\label{ch:protocol}

The protocol is written before any experiment runs, so that the results
cannot shape the questions. It specifies the metrics, the baselines, the
ablations, the robustness tests, the statistical treatment, the latency
measurement and the compute budget. Each choice answers a demand that
reviewers at the target journals make routinely (Section~\ref{sec:reviewers}).

\section{Metrics}
\label{sec:metrics}

\paragraph{Still images}
COCO AP, AP\textsubscript{50} and AP\textsubscript{75} are computed with
\texttt{pycocotools} or the faster compatible
\texttt{faster-coco-eval}~\citep{lin2014coco,fastercocoeval_repo}. Ultralytics
reports only AP\textsubscript{50} and AP, so predictions are exported as COCO
JSON and scored externally. Small-object performance is reported in three
ways:
\begin{itemize}
  \item the COCO size bands;
  \item the AI-TOD bands of 2--8, 8--16 and 16--32 pixels;
  \item SAFit-based AP on RGBT-Tiny~\citep{ying2025rgbttiny}.
\end{itemize}
If KAIST is used, the log-average miss rate is also reported.

\paragraph{Video and tracking}
Video detection is scored per frame over whole sequences. Tracking uses
TrackEval~\citep{luiten2021hota,trackeval_repo}, which reports HOTA with its
detection and association components, IDF1 and MOTA. Temporal stability is
measured as \emph{flicker}. Flicker is the fraction of ground-truth track
frames in which a detected object is lost and found again within three
frames.

\paragraph{Streaming robustness}
Each perturbation of Section~\ref{sec:robustness} is applied at several
severities. The result is a degradation curve of AP against severity. Its
normalised area (AUDC) summarises robustness in one number. Two event
metrics describe behaviour around outages.
\begin{itemize}
  \item \emph{Hold recall} is the recall, during an outage, of objects that
        were visible before it began.
  \item \emph{Recovery time} $R_{95}$ is the number of frames after a sensor
        returns until per-frame recall reaches 95\% of the recall on the
        same frames without the outage.
\end{itemize}

\paragraph{Efficiency}
The efficiency metrics are parameters and GFLOPs at $640\times512$, and
end-to-end latency at batch one (Section~\ref{sec:latency}). They also
include throughput, peak GPU memory, and memory as a function of sequence
length. For \methodS{}, the memory curve should be flat.

\section{Baselines}
\label{sec:baselines}

Every baseline is retrained under the same schedule, input size and
evaluation code. None are copied from papers. Where official code exists it
is used; otherwise the method is reimplemented and the differences are
stated. Table~\ref{tab:baselines} lists the four groups.

\begin{table}[htbp]
\centering
\caption[Baselines by group]{Baselines by group. The code column records whether an
official implementation is expected to be available; reimplementations are
labelled as such in the paper.}
\label{tab:baselines}
\small
\begin{tabularx}{\textwidth}{@{}P{2.55cm}LP{2.2cm}@{}}
\hd{Group} & \hd{Methods} & \hd{Code} \\ \gap
single modality, per frame & YOLO26-s on visible; YOLO26-s on thermal; RF-DETR-S on visible as a DETR reference~\citep{robinson2025rfdetr} & official \\ \sgap
static fusion & two-stream YOLO26 with concatenation and addition; CFT-style attention fusion~\citep{qingyun2021cft}; Fusion-Mamba~\citep{dong2025fusionmamba}; COMO~\citep{como2024}; WaveMamba~\citep{wavemamba2025}; image-mode \method{} & mixed; reimplement where missing \\ \sgap
temporal & sliding-window feature aggregation with $k\in\{3,5,9\}$; convolutional GRU~\citep{ballas2016convgru}; windowed temporal attention; three-frame MambaST-style fuser~\citep{mambast2024}; bidirectional temporal SSM as an offline upper bound; MSENet or MSGNet if released~\citep{msenet2024,msgnet2025} & mostly reimplemented \\ \sgap
tracking & ByteTrack, BoT-SORT and TrackTrack on top of each detector~\citep{zhang2022bytetrack,aharon2022botsort,shim2025tracktrack} & official, in Ultralytics \\
\end{tabularx}
\end{table}

The temporal baselines are the critical group. Each receives the same
per-frame fused features $Z_t$ as the TSM and replaces only the temporal
module. The comparison therefore isolates the memory mechanism. This is the
control that MambaOut-style critiques demand~\citep{yu2025mambaout}. It is
complemented by a \emph{gated-convolution control}: the TSM with its
state-space recurrence replaced by a gated convolution of matched cost.

\section{Ablations}
\label{sec:ablations}

Table~\ref{tab:ablations} lists the ablations, each tied to a design
decision of Chapter~\ref{ch:method}. To fit the compute budget, every ablation
first runs with one seed. Those that change the conclusion of a hypothesis
are repeated with three seeds.

\begin{table}[htbp]
\centering
\caption[The ablation plan]{The ablation plan. The default setting is listed
first in each row.}
\label{tab:ablations}
\small
\begin{tabularx}{\textwidth}{@{}P{0.75cm}P{3.6cm}LP{0.95cm}@{}}
\hd{} & \hd{Component} & \hd{Variants} & \hd{Tests} \\ \gap
A1  & stride-4 head & on; off & P1 \\
A2  & early capacity & UAVDet-style; YOLO26 default; intermediate & P1 \\
A3  & fusion operator & CMFM with gated step; CMFM with output gating; attention; concatenation; addition & H3 \\
A4  & token order & interleaved; separate scans then merge & H3 \\
A5  & scan directions & four; two & cost \\
A6  & offset alignment & on; off & H3 \\
A7  & TSM placement & $P_3$--$P_5$; $P_3$ only; $P_3$--$P_4$; including $P_2$ & H1 \\
A8  & state size $N$ & 8; 4; 16 & H1 \\
A9  & motion compensation & classical homography; learned; none & H1 \\
A10 & interval scaling $\tau_t/\tau_0$ & on; off & H4 \\
A11 & reliability supervision & on; off & H3 \\
A12 & clip training & $L=16$ with carry 0.5; $L=8$; no carry & H1 \\
A13 & reset & hard; soft; none & H1 \\
A14 & memory mechanism & TSM; gated convolution; ConvGRU; temporal attention & H2 \\
A15 & backbones & separate; shared from stage 3 & cost \\
\end{tabularx}
\end{table}

\section{Robustness protocol}
\label{sec:robustness}

No benchmark protocol exists for temporal sensor failure in visible--thermal
detection (Section~\ref{sec:directions}). This section defines one, and it
is itself a contribution of the paper (Table~\ref{tab:robustness}). All
perturbations are applied at test time to the RGBT-Tiny test sequences with
fixed random seeds, so that every method sees identical corrupted streams.
Real misalignment is tested on UVT-VOD2024 and on the shifted subset of
DroneVehicle~\citep{msenet2024,cbcswca2025}.

\begin{table}[htbp]
\centering
\caption[The streaming robustness protocol]{The streaming robustness protocol.
Severities are applied separately unless combined in the last row.}
\label{tab:robustness}
\small
\begin{tabularx}{\textwidth}{@{}P{3.2cm}LP{3.0cm}@{}}
\hd{Perturbation} & \hd{Parameters and severities} & \hd{Metrics} \\ \gap
independent dropout & each modality dropped per frame with probability 0.1, 0.3, 0.5 & AP, AUDC \\
bursty dropout & Markov outages with mean length 5, 15, 45 frames & hold recall, $R_{95}$ \\
blackout & both sensors absent for 5 or 15 frames & hold recall, $R_{95}$ \\
misalignment & thermal translation 0, 4, 8, 16\,px; rotation 0, 1, 3 degrees; scale 0, 2, 5\% & AP, AUDC \\
drift & offset growing linearly to the maximum over 300 frames & AP over time \\
desynchronisation & thermal delayed by 0, 1, 2 frames & AP \\
frame rate & strides 2 and 4 (7.5 and 3.75 FPS) and random drops of 10\% and 30\% & AP, flicker \\
sensor degradation & thermal gain drift and noise; visible low light and synthetic fog & AP \\
combined & bursty thermal dropout with 8\,px misalignment and 10\% frame drops & all \\
\end{tabularx}
\end{table}

\section{Statistical practice}
\label{sec:stats}

\begin{itemize}
  \item Every main configuration is trained with three seeds and reported as
        mean and standard deviation.
  \item Frames within a sequence are not independent, so confidence intervals
        are computed by bootstrap over \emph{sequences}, with 10\,000
        resamples.
  \item Comparisons between two methods use the paired bootstrap of their
        difference on the same resampled sequences. A difference is reported
        as significant when its 95\% interval excludes zero.
  \item The five hypotheses form a family. Holm's correction is applied to
        their tests.
  \item Effect sizes are reported in AP points, not only as significance.
\end{itemize}

\section{Latency and efficiency measurement}
\label{sec:latency}

Latency is measured by one script for every model, on the same machine, with
the software versions recorded. The protocol follows the cautions of
RF-DETR~\citep{robinson2025rfdetr}.
\begin{itemize}
  \item The input is batch one at $640\times512$, with FP16 TensorRT engines.
  \item Warm-up runs 50 iterations. Measurement runs 500 iterations, with a
        200\,ms pause between them to avoid thermal throttling. The median
        and the 95th percentile are reported.
  \item The timed region includes pre-processing, the network, NMS where the
        head needs it, motion estimation and the TSM update.
  \item Accuracy is measured on the exported engine, not on the PyTorch
        model, so that any loss in export is visible.
  \item On Jetson Orin, the power mode and clock settings are stated.
\end{itemize}

\section{Compute budget}
\label{sec:compute}

The work runs on two resources (Table~\ref{tab:hardware}). A local laptop
GPU handles development, testing and live demonstrations. Kaggle notebooks
with T4 GPUs handle training. Kaggle runs Linux, so the \texttt{mamba\_ssm}
kernels can be compiled there. The compiled wheel is stored as a private
Kaggle dataset so that later sessions install it in seconds. The T4 has
compute capability 7.5 and no bfloat16. Kernel support on it must be
verified in Notebook~00, with the pure-PyTorch scan as the
fallback~\citep{mambapy_repo}.

\begin{table}[htbp]
\centering
\caption[Compute resources and their roles]{Compute resources and their roles.
Kaggle limits change over time and must be checked when the work starts.}
\label{tab:hardware}
\small
\begin{tabularx}{\textwidth}{@{}P{3.4cm}LL@{}}
\hd{Resource} & \hd{Specification} & \hd{Role} \\ \gap
local laptop & RTX 4050 Laptop GPU, 6\,GB, Windows 11; WSL2 Ubuntu for Linux tools & unit tests, debugging at small scale, visualisation, live demonstration, laptop latency \\ \sgap
Kaggle notebooks & two T4 GPUs with 16\,GB each, or one P100; sessions of up to 12\,h; weekly GPU quota of about 30\,h\unv & all training and evaluation runs; distributed data parallel over two T4s \\ \sgap
optional & university GPU cluster, Colab, a borrowed Jetson Orin & contingency for quota; edge latency for H5 \\
\end{tabularx}
\end{table}

Table~\ref{tab:budget} estimates the training budget in T4 GPU-hours. The
estimates rest on three assumptions, which Notebook~04 replaces with
measurements:
\begin{itemize}
  \item the small (s) model scale;
  \item native $640\times512$ input with COCO-pretrained backbones;
  \item RGBT-Tiny training frames subsampled by three, which loses little
        because neighbouring frames at 15 FPS are nearly identical.
\end{itemize}
About 440 T4-hours fit into roughly fifteen weeks of Kaggle quota, which is
compatible with the timeline in Chapter~\ref{ch:plan}. Running sessions
continuously and using both T4s per session reduces wall-clock time.

\begin{table}[htbp]
\centering
\caption[Estimated training budget in T4 GPU-hours]{Estimated training budget
in T4 GPU-hours. All figures are planning estimates to be replaced by
measurements.}
\label{tab:budget}
\small
\begin{tabularx}{\textwidth}{@{}LR{2.0cm}R{1.4cm}R{1.7cm}R{1.5cm}@{}}
\hd{Experiment group} & \hd{Configurations} & \hd{Seeds} & \hd{Hours each} & \hd{Total} \\ \gap
Phase 1 baselines and \method{} on RGBT-Tiny & 7 & 3 & 4 & 84 \\
image ablations A1--A6, A15 & 10 & 1 & 4 & 40 \\
confirmation of decisive image ablations & 3 & 2 & 4 & 24 \\
Stage B: \methodS{} and temporal baselines & 6 & 3 & 6 & 108 \\
temporal ablations A7--A14 & 10 & 1 & 6 & 60 \\
other datasets: VisDrone, DroneVehicle, VT-VOD50, UVT-VOD2024, KAIST & 15 & 1--3 & 4 & 90 \\
robustness evaluation, tracking, export & & & & 30 \\ \gap
\textit{total} & & & & \textit{436} \\
\end{tabularx}
\end{table}

Kaggle sessions end after at most twelve hours. Every training script
therefore writes a checkpoint at the end of each epoch and resumes
automatically from the latest one. Runs are sized to finish within a
session wherever possible.
